% =====================================================================
%  AEGIS v2 — presentation (ARGUS Open Innovation Challenge 2026)
%  Compile with XeLaTeX:   xelatex AEGIS_presentation_v2.tex
%
%  Fonts (Carlito, Liberation Mono) are bundled in ./fonts and all
%  pictures / screenshot shadows are in ./images.
%
%  Coordinate system: every slide is the original 1280 x 720 px canvas
%  (1px = 0.75bp, so a page is 960 x 540 bp), origin at the top-left
%  corner, y growing downwards.  Every box, picture and line of text is
%  given in those pixels, exactly where the original slide has it.
%
%  Text:    \T(x,y){...}  puts a line of text, baseline origin at (x,y)
%           \R{size px}{RRGGBB}{text}   Carlito Regular
%           \B{size px}{RRGGBB}{text}   Carlito Bold
%           \M{size px}{RRGGBB}{text}   Liberation Mono
%           \LS{percent}{text}          CSS letter-spacing (in % of em)
%  Shapes:  \Fill{RRGGBB}(x0,y0)(x1,y1)             filled rectangle
%           \RoundFill{RRGGBB}{r}{nw,ne}(x0,y0)(x1,y1) top corners rounded
%           \Frame{RRGGBB}{w}(x0,y0)(x1,y1)          stroked rectangle
%           \Img{file}(x0,y0)(x1,y1)                 picture in a box
%           \Shadow{file}{opacity}(x0,y0)(x1,y1)[box] CSS box-shadow
% =====================================================================
\documentclass{article}
\usepackage[paperwidth=960bp,paperheight=540bp,margin=0pt]{geometry}
\usepackage{fontspec}
\usepackage{xcolor}
\usepackage{graphicx}
\usepackage{eso-pic}
% keep 5 decimals in the PDF so nothing is rounded to the pixel grid
\special{dvipdfmx:config d 5}

\defaultfontfeatures{Path=fonts/}
\newfontfamily\carlito{Carlito-Regular.ttf}
\newfontfamily\carlitobold{Carlito-Bold.ttf}
\newfontfamily\libmono{LiberationMono-Regular.ttf}
\def\fs#1{\fontsize{\fpeval{#1*0.75}bp}{\fpeval{#1*0.75}bp}\selectfont}
\newcommand\R[3]{{\carlito\fs{#1}\color[HTML]{#2}\words{#3}}}
\newcommand\B[3]{{\carlitobold\fs{#1}\color[HTML]{#2}\words{#3}}}
% Chrome lays text out at the font size rounded to 1/64 px (12.6px ->
% 12.59375px); for the monospaced listing that shifts every letter a
% little, so \M adds the same tiny correction after each character.
\newcommand\M[3]{{\libmono\fs{#1}\color[HTML]{#2}%
  \edef\monofix{\fpeval{(round(#1*64)/64-(#1))*1229/2048*0.75}bp}\chars{#3}}}
% CSS letter-spacing, in percent of the font size
\newcommand\LS[2]{{\addfontfeature{LetterSpace=#1}%
  \spaceskip=\dimexpr\fontdimen2\font+#1\dimexpr0.01em\relax\relax\words{#2}}}
% \W{word}: a word shaped with a ligature, typeset in one piece
\newcommand\W[1]{#1}
% \words typesets its argument word by word, and each word glyph by
% glyph: every glyph is put at the exact pen position TeX computed for it
% (kerning included), and every other glyph sits 0.0002bp higher.  That is
% invisible, but it makes the PDF driver re-anchor each glyph instead of
% letting the rounding of its 1/1000-em glyph widths add up along a line,
% so the letters land on the same sub-pixel positions as in the original.
\ExplSyntaxOn
\int_new:N \g__aegis_word_int
\box_new:N \l__aegis_w_box
\box_new:N \l__aegis_p_box
\box_new:N \l__aegis_g_box
\dim_new:N \l__aegis_x_dim
\tl_new:N \l__aegis_prefix_tl
\cs_new_protected:Npn \__aegis_word:n #1
  {
    \int_gincr:N \g__aegis_word_int
    \raise \int_if_odd:nTF { \g__aegis_word_int } { 0.0002bp } { 0bp } \hbox {#1}
  }
\cs_new_protected:Npn \__aegis_glyphs:n #1
  {
    \tl_if_in:nnTF {#1} { \LS } {#1}
      {
      \tl_if_in:nnTF {#1} { \W } { \__aegis_word:n {#1} }
      {
        \tl_clear:N \l__aegis_prefix_tl
        \hbox_set:Nn \l__aegis_w_box {#1}
        \hbox_to_wd:nn { \box_wd:N \l__aegis_w_box }
          {
            \tl_map_inline:nn {#1}
              {
                \tl_put_right:Nn \l__aegis_prefix_tl {##1}
                \hbox_set:Nn \l__aegis_p_box { \l__aegis_prefix_tl }
                \hbox_set:Nn \l__aegis_g_box {##1}
                \dim_set:Nn \l__aegis_x_dim
                  { \box_wd:N \l__aegis_p_box - \box_wd:N \l__aegis_g_box }
                \hbox_overlap_right:n
                  { \skip_horizontal:N \l__aegis_x_dim \__aegis_word:n {##1} }
              }
            \hfil
          }
      }
      }
  }
\cs_new_protected:Npn \words #1
  {
    \seq_set_split:Nnn \l_tmpa_seq { ~ } {#1}
    \seq_set_map:NNn \l_tmpb_seq \l_tmpa_seq { \__aegis_glyphs:n {##1} }
    \seq_use:Nn \l_tmpb_seq { ~ }
  }
% same, letter by letter (used for the monospaced code listing)
\cs_new_protected:Npn \chars #1
  {
    \tl_set:Nn \l_tmpa_tl {#1}
    \tl_replace_all:Nnn \l_tmpa_tl { ~ } { \  }
    \tl_map_inline:Nn \l_tmpa_tl { \__aegis_word:n {##1} \kern \monofix }
  }
\ExplSyntaxOff

% ------------------------------------------------------------- machinery
% px -> sp for \put (1px = 49336.32sp) and px -> bp for PDF code; done
% with exact decimal arithmetic so positions are not rounded.
\def\spx#1{\fpeval{round((#1)*49336.32)}}
\def\spy#1{\fpeval{round(-(#1)*49336.32)}}
\def\bp#1{\fpeval{round((#1)*0.75,5)}}
\def\nbp#1{\fpeval{round(-(#1)*0.75,5)}}
\def\pdf#1{\put(0,0){\special{pdf:literal #1}}}
% RRGGBB -> "r g b"
\def\rgb#1{\convertcolorspec{HTML}{#1}{rgb}\tmpc\expandafter\rgbsplit\tmpc\relax}
\def\rgbsplit#1,#2,#3\relax{\def\tmpc{#1 #2 #3}}

\def\T(#1,#2)#3{\put(\spx{#1},\spy{#2}){#3}}
\def\Fill#1(#2,#3)(#4,#5){\rgb{#1}%
  \pdf{\tmpc\space rg \bp{#2} \nbp{#5} \bp{#4-#2} \bp{#5-#3} re f}}
\def\RoundFill#1#2#3(#4,#5)(#6,#7){\rgb{#1}%
  \pdf{\tmpc\space rg \bp{#4} \nbp{#7} m \bp{#4} \nbp{#5+#2} l
       \bp{#4} \nbp{#5+0.4477*#2} \bp{#4+0.4477*#2} \nbp{#5} \bp{#4+#2} \nbp{#5} c
       \bp{#6-#2} \nbp{#5} l
       \bp{#6-0.4477*#2} \nbp{#5} \bp{#6} \nbp{#5+0.4477*#2} \bp{#6} \nbp{#5+#2} c
       \bp{#6} \nbp{#7} l h f}}
\def\Frame#1#2(#3,#4)(#5,#6){\rgb{#1}%
  \pdf{\tmpc\space RG \bp{#2} w \bp{#3} \nbp{#6} \bp{#5-#3} \bp{#6-#4} re S}}
\def\Img#1(#2,#3)(#4,#5){\put(\spx{#2},\spy{#5}){%
  \special{pdf:image width \bp{#4-#2}bp height \bp{#5-#3}bp (#1)}}}
\def\Shadow#1#2(#3,#4)(#5,#6)[#7,#8,#9]{%
  \def\shfile{#1}\def\shop{#2}\def\shxa{#3}\def\shya{#4}\def\shxb{#5}\def\shyb{#6}%
  \ShadowHole{#7}{#8}#9\relax}
\def\ShadowHole#1#2#3,#4\relax{%
  \special{pdf:put @resources << /ExtGState << /AegisShadow << /ca \shop\space >> >> >>}%
  \pdf{q 0 -540 960 540 re \bp{#1} \nbp{#4} \bp{#3-#1} \bp{#4-#2} re W* n /AegisShadow gs}%
  \Img{\shfile}(\shxa,\shya)(\shxb,\shyb)%
  \pdf{Q}}

% one slide = one page; drawing origin = top-left corner of the page
\NewDocumentEnvironment{slide}{+b}
  {\AddToShipoutPictureFG*{\AtPageUpperLeft{\setlength\unitlength{1sp}#1}}\null\newpage}{}

\pagestyle{empty}
\frenchspacing % browsers do not add extra space after punctuation
\begin{document}
% ================================================= slide 1 — Τίτλος
\begin{slide}
\Fill{195A49}(0,0)(1280,720)
\Img{images/slide01_1.jpg}(40,30)(144,143)
\Img{images/slide01_2.png}(40,510)(168,638)
\Shadow{images/slide01_shadow1.png}{0.1608}(605.12,161.28)(1246.72,531.84)[616,170,1236,519]
\Img{images/slide01_3.jpg}(616,170)(1236,519)
% text
\T(40,222){\B{76}{FFFFFF}{\LS{1}{AEGIS}}}
\T(40,276){\R{31}{FFFFFF}{Ο επισκέπτης φύλακας του μνημείου}}
\T(40,366){\R{21}{FFFFFF}{ARGUS Open \W{Innovation} Challenge 2026}}
\T(40,394){\R{21}{FFFFFF}{Track A: Πολιτιστική Κληρονομιά και Επιστήμη των Πολιτών}}
\T(196,541){\R{20}{FFFFFF}{Παίξτε το τώρα, σε υπολογιστή}}
\T(196,569){\R{20}{FFFFFF}{ή κινητό:}}
\T(196,597){\M{18}{FFFFFF}{https://aloualou56.github.io/}}
\T(196,626){\M{18}{FFFFFF}{argus\_trackA/}}
\end{slide}

% ================================================= slide 2 — Το πρόβλημα
\begin{slide}
\Fill{FBFCFA}(0,0)(1280,720)
\Fill{0B5D4B}(58,109)(1222,111)
\RoundFill{EDF2EB}{6}{nw,ne}(58,217)(662,461)
\RoundFill{0B5D4B}{6}{nw,ne}(58,217)(662,260)
\RoundFill{EDF2EB}{6}{nw,ne}(688,217)(1222,461)
\RoundFill{0B5D4B}{6}{nw,ne}(688,217)(1222,260)
% text
\T(58,77){\B{40}{0B5D4B}{Το πρόβλημα: η δουλειά του ARGUS δεν φτάνει στον κόσμο}}
\T(58,160){\R{22}{14211B}{Το ARGUS παρακολουθεί πέντε μνημεία της Ευρώπης. Καλείται να υπολογίζει καθημερινά τους κινδύνους που τα απειλούν και}}
\T(58,191){\R{22}{14211B}{να σημαίνει συναγερμό όταν κάτι ξεφεύγει. Όταν όμως εξετάσαμε τα δεδομένα του, ήρθαμε αντιμέτωποι με το εξής:}}
\T(70,247){\R{24}{FFFFFF}{Ένας συναγερμός, όπως τον γράφει το ARGUS}}
\T(70,279){\M{12.6}{14211B}{\{"type": "DL",}}
\T(70,296){\M{12.6}{14211B}{\ "score": 0.026571929454803467,}}
\T(70,312){\M{12.6}{14211B}{\ "threshold": 0.020811407233135213,}}
\T(70,329){\M{12.6}{14211B}{\ "is\_alarm": true,}}
\T(70,346){\M{12.6}{14211B}{\ "confidence\_score": 0.5004581977196843,}}
\T(70,362){\M{12.6}{14211B}{\ "top\_feature\_contributions": \{}}
\T(70,379){\M{12.6}{14211B}{\ \ "149368\_h\_masks\_in\_room215\_2\_soil\_moisture": 0.4397,}}
\T(70,396){\M{12.6}{14211B}{\ \ "151836\_h\_masks\_in\_room214\_hum": 0.2579,}}
\T(70,412){\M{12.6}{14211B}{\ \ "168439\_h\_masks\_in\_room215\_lux": 0.0625, ...\},}}
\T(70,429){\M{12.6}{14211B}{\ "site": "delos",}}
\T(70,445){\M{12.6}{14211B}{\ "timestamp": "2026-08-31T20:30:00"\}}}
\T(699.719,247){\R{24}{FFFFFF}{Ο ίδιος συναγερμός, όπως εμφανίζεται στο AEGIS}}
\T(699.719,288){\R{20}{14211B}{«Στις 31 Αυγούστου 2026, το σύστημα τεχνητής νοημοσύνης}}
\T(699.719,316){\R{20}{14211B}{του ARGUS πρόσεξε κάτι ασυνήθιστο στην υγρασία του}}
\T(699.719,343){\R{20}{14211B}{εδάφους κάτω από ένα ψηφιδωτό.»}}
\T(699.719,381){\R{20}{14211B}{«Αυτός ο αισθητήρας εξηγεί το 44\% του συναγερμού,}}
\T(699.719,409){\R{20}{14211B}{περισσότερο από κάθε άλλον.»}}
\T(58,500){\R{20}{14211B}{Με βάση αυτά, κρίναμε ότι ο πιο προσιτός τρόπος για να κατανοήσει ο απλός πολίτης το έργο του ARGUS είναι η }\B{20}{14211B}{βιωματική μάθηση μέσα}}
\T(58,528){\B{20}{14211B}{από ένα σοβαρό παιχνίδι ρόλων}\R{20}{14211B}{. Το παιχνίδι βασίζεται σε αυθεντικά δεδομένα παρακολούθησης μνημείων, ακολουθεί τον κύκλο μάθησης}}
\T(58,557){\R{20}{14211B}{του Kolb και επεκτείνεται στον πραγματικό κόσμο μέσω της επιστήμης των πολιτών.}}
\T(58,588){\R{14}{5F6863}{Αριστερά: η εγγραφή 123 του αρχείου συμβάντων του ARGUS, με λιγότερα πεδία. Δεξιά: η οθόνη του συναγερμού στο παιχνίδι.}}
\T(1157.297,701){\R{13}{88928C}{\LS{2}{AEGIS}}}
\T(1201.766,701){\R{13}{88928C}{\LS{2}{2 / 10}}}
\end{slide}

% ================================================= slide 3 — Γιατί υπάρχει
\begin{slide}
\Fill{FBFCFA}(0,0)(1280,720)
\Fill{0B5D4B}(58,109)(1222,111)
\Shadow{images/slide03_shadow1.png}{0.1608}(691.2,132.16)(1232.96,435.84)[702,141,1222,423]
\Frame{D9DED9}{1}(702.5,141.5)(1221.5,422.5)
\Img{images/slide03_1.jpg}(703,142)(1221,422)
\Fill{A84300}(58,150)(69,161)
\Fill{A84300}(58,255)(69,266)
\Fill{A84300}(58,299)(69,310)
\Fill{A84300}(58,344)(69,355)
% text
\T(58,77){\B{40}{0B5D4B}{Γιατί υπάρχει}}
\T(702,454){\R{17}{14211B}{Μέσα από το παιχνίδι, ο χρόνος τέμνεται στα δύο: το «Σήμερα»}}
\T(702,479){\R{17}{14211B}{αναμετράται με τα «Χρόνια μετά, αν κανείς δεν προσέξει». Για να φανεί τι}}
\T(702,503){\R{17}{14211B}{θα συμβεί αν δεν προλάβουμε τη φθορά, χρησιμοποιούνται ζωγραφιές}}
\T(702,528){\R{17}{14211B}{και όχι φωτογραφίες του μνημείου, και το παιχνίδι το δηλώνει ξεκάθαρα.}}
\T(84,163){\R{22}{14211B}{Η φθορά σε ένα μνημείο είναι }\B{22}{14211B}{αργή και αθόρυβη}\R{22}{14211B}{. Η υγρασία}}
\T(84,193){\R{22}{14211B}{αποσαθρώνει τους τοίχους, το αλάτι της θάλασσας διαβρώνει}}
\T(84,224){\R{22}{14211B}{την πέτρα, το ψηφιδωτό χάνει τις ψηφίδες του.}}
\T(84,268){\R{22}{14211B}{Όταν η ζημιά γίνει ορατή, }\B{22}{14211B}{είναι πλέον αργά}\R{22}{14211B}{.}}
\T(84,312){\R{22}{14211B}{Η εξειδίκευση }\B{22}{14211B}{περιχαρακώνει τη γνώση}\R{22}{14211B}{.}}
\T(84,357){\R{22}{14211B}{Έτσι, ο απλός πολίτης βλέπει το αποτέλεσμα, όχι τη δουλειά}}
\T(84,387){\R{22}{14211B}{πίσω από αυτό. Κατά συνέπεια, δεν πιστεύει ότι μπορεί να}}
\T(84,417){\R{22}{14211B}{συμβάλει κι ο ίδιος.}}
\T(1157.297,701){\R{13}{88928C}{\LS{2}{AEGIS}}}
\T(1201.766,701){\R{13}{88928C}{\LS{2}{3 / 10}}}
\end{slide}

% ================================================= slide 4 — Η ιδέα μας
\begin{slide}
\Fill{FBFCFA}(0,0)(1280,720)
\Fill{0B5D4B}(58,109)(1222,111)
\Shadow{images/slide04_shadow1.png}{0.1608}(651.2,132.16)(1232.96,469.76)[662,141,1222,457]
\Frame{D9DED9}{1}(662.5,141.5)(1221.5,456.5)
\Img{images/slide04_1.jpg}(663,142)(1221,456)
\Fill{A84300}(58,150)(69,161)
\Fill{A84300}(58,225)(69,236)
\Fill{A84300}(58,299)(69,310)
\Fill{A84300}(58,374)(69,385)
% text
\T(58,77){\B{40}{0B5D4B}{Η ιδέα μας: φύλακας για μια μέρα}}
\T(84,163){\R{22}{14211B}{Σύμφωνα με τα παραπάνω, είναι φανερό ότι η δουλειά του}}
\T(84,193){\R{22}{14211B}{ARGUS }\B{22}{14211B}{δεν γίνεται αντιληπτή μέσα από λέξεις}\R{22}{14211B}{.}}
\T(84,238){\R{22}{14211B}{Στο AEGIS ο καθένας γίνεται }\B{22}{14211B}{φύλακας σε ένα από τα πέντε}}
\T(84,268){\B{22}{14211B}{μνημεία}\R{22}{14211B}{\ του ARGUS, για έως και 7 μέρες.}}
\T(84,312){\R{22}{14211B}{Ένας επινοημένος φύλακας, π.χ. ο Μανώλης στη Δήλο,}}
\T(84,343){\R{22}{14211B}{παίρνει ρεπό και του }\B{22}{14211B}{παραδίδει τη βάρδια}\R{22}{14211B}{.}}
\T(84,387){\R{22}{14211B}{Το παιχνίδι επισημαίνει ότι ο φύλακας μπορεί να είναι}}
\T(84,417){\R{22}{14211B}{φανταστικός, αλλά οι κίνδυνοι και οι συναγερμοί είναι}}
\T(84,448){\B{22}{14211B}{πέρα για πέρα αληθινοί}\R{22}{14211B}{.}}
\T(1157.297,701){\R{13}{88928C}{\LS{2}{AEGIS}}}
\T(1201.766,701){\R{13}{88928C}{\LS{2}{4 / 10}}}
\end{slide}

% ================================================= slide 5 — Πώς παίζεται (1/2)
\begin{slide}
\Fill{FBFCFA}(0,0)(1280,720)
\Fill{0B5D4B}(58,109)(1222,111)
\Shadow{images/slide05_shadow1.png}{0.1608}(47.04,304)(628.8,641.92)[58,313,618,629]
\Frame{D9DED9}{1}(58.5,313.5)(617.5,628.5)
\Img{images/slide05_1.jpg}(59,314)(617,628)
\Shadow{images/slide05_shadow2.png}{0.1608}(651.2,304)(1232.96,641.92)[662,313,1222,629]
\Frame{D9DED9}{1}(662.5,313.5)(1221.5,628.5)
\Img{images/slide05_2.jpg}(663,314)(1221,628)
% text
\T(58,77){\B{40}{0B5D4B}{Πώς παίζεται: το δίκτυό σου, ένας πραγματικός συναγερμός}}
\T(58,162){\B{20.5}{A84300}{1. Φτιάξε το δίκτυό σου.}\R{20.5}{14211B}{\ Ο παίκτης καλείται, με 12 έως 16}}
\T(58,190){\R{20.5}{14211B}{νομίσματα ανάλογα με το επίπεδο κινδύνου, να διαλέξει και να}}
\T(58,219){\R{20.5}{14211B}{τοποθετήσει πάνω στο μνημείο τους κατάλληλους αισθητήρες για}}
\T(58,248){\R{20.5}{14211B}{τρεις πραγματικούς κινδύνους.}}
\T(662,162){\B{20.5}{A84300}{2. Ξαναζήσε έναν πραγματικό συναγερμό.}\R{20.5}{14211B}{\ Ο παίκτης έρχεται}}
\T(662,190){\R{20.5}{14211B}{αντιμέτωπος με αληθινούς συναγερμούς, όπως αυτόν που σήμανε}}
\T(662,219){\R{20.5}{14211B}{το ARGUS στις 31 Αυγούστου 2026, ακριβώς στο ίδιο σημείο όπου}}
\T(662,248){\R{20.5}{14211B}{χτύπησε, πάνω στο 3D μοντέλο του ίδιου του ARGUS, όπου}}
\T(662,277){\R{20.5}{14211B}{υπάρχει.}}
\T(1157.297,701){\R{13}{88928C}{\LS{2}{AEGIS}}}
\T(1201.766,701){\R{13}{88928C}{\LS{2}{5 / 10}}}
\end{slide}

% ================================================= slide 6 — Πώς παίζεται (2/2)
\begin{slide}
\Fill{FBFCFA}(0,0)(1280,720)
\Fill{0B5D4B}(58,109)(1222,111)
\Shadow{images/slide06_shadow1.png}{0.1608}(47.04,304)(628.8,641.92)[58,313,618,629]
\Frame{D9DED9}{1}(58.5,313.5)(617.5,628.5)
\Img{images/slide06_1.jpg}(59,314)(617,628)
\Shadow{images/slide06_shadow2.png}{0.1608}(651.2,304)(1232.96,641.92)[662,313,1222,629]
\Frame{D9DED9}{1}(662.5,313.5)(1221.5,628.5)
\Img{images/slide06_2.jpg}(663,314)(1221,628)
% text
\T(58,77){\B{40}{0B5D4B}{Πώς παίζεται: θα το είχες προσέξει; τι έβαλε το ARGUS}}
\T(58,162){\B{20.5}{A84300}{3. Θα το είχε πιάσει το δίκτυό σου;}\R{20.5}{14211B}{\ Το ARGUS σού δείχνει ποιος}}
\T(58,190){\R{20.5}{14211B}{αισθητήρας χτύπησε και πόσες φορές.}}
\T(662,162){\B{20.5}{A84300}{4. Τι έβαλε το ARGUS, και πού.}\R{20.5}{14211B}{\ Τα 30 όργανά του στη Δήλο, στις}}
\T(662,190){\R{20.5}{14211B}{πραγματικές τους θέσεις, δίπλα δίπλα με το δίκτυο που έστησε ο}}
\T(662,219){\R{20.5}{14211B}{παίκτης. Το παιχνίδι εξηγεί ότι το δίκτυο του ARGUS θα κόστιζε 52}}
\T(662,248){\R{20.5}{14211B}{νομίσματα, αλλά ακόμα κι αυτό έπρεπε να επιλέξει προσεκτικά τι}}
\T(662,277){\R{20.5}{14211B}{θα τοποθετήσει και πού.}}
\T(1157.297,701){\R{13}{88928C}{\LS{2}{AEGIS}}}
\T(1201.766,701){\R{13}{88928C}{\LS{2}{6 / 10}}}
\end{slide}

% ================================================= slide 7 — Αν κανείς δεν το πρόσεχε
\begin{slide}
\Fill{FBFCFA}(0,0)(1280,720)
\Fill{0B5D4B}(58,109)(1222,111)
\Shadow{images/slide07_shadow1.png}{0.1608}(47.04,128)(708.8,510.72)[58,137,698,498]
\Frame{D9DED9}{1}(58.5,137.5)(697.5,497.5)
\Img{images/slide07_1.jpg}(59,138)(697,497)
\Shadow{images/slide07_shadow2.png}{0.1608}(721.28,233.28)(1102.72,593.92)[732,242,1092,581]
\Frame{D9DED9}{1}(732.5,242.5)(1091.5,580.5)
\Img{images/slide07_2.jpg}(733,243)(1091,580)
% text
\T(58,77){\B{40}{0B5D4B}{Αν κανείς δεν το πρόσεχε}}
\T(58,529){\R{18.5}{14211B}{Έως 50 χρόνια χωρίς φροντίδα, πάνω στο 3D μοντέλο του ARGUS για την Οικία με τα}}
\T(58,555){\R{18.5}{14211B}{Προσωπεία: η καταστροφή γίνεται βίωμα. Το νερό ανεβαίνει από το χώμα, το αλάτι}}
\T(58,581){\R{18.5}{14211B}{που φέρνει κρυσταλλώνει, και οι ψηφίδες του ψηφιδωτού ξεκολλούν και χάνονται.}}
\T(58,607){\R{18.5}{14211B}{Η ζημιά είναι εικονογράφηση, και η οθόνη το γράφει.}}
\T(732,158){\B{21}{14211B}{Ο φύλακας επιστρέφει}\R{21}{14211B}{\ και εξηγεί γιατί, όταν}}
\T(732,188){\R{21}{14211B}{επισκεπτόμαστε ένα μνημείο, είναι σημαντικό να}}
\T(732,217){\R{21}{14211B}{ακολουθούμε τους κανόνες του.}}
\T(1157.297,701){\R{13}{88928C}{\LS{2}{AEGIS}}}
\T(1201.766,701){\R{13}{88928C}{\LS{2}{7 / 10}}}
\end{slide}

% ================================================= slide 8 — Τι πήραμε από το ARGUS
\begin{slide}
\Fill{FBFCFA}(0,0)(1280,720)
\Fill{0B5D4B}(58,109)(1222,111)
\Fill{14211B}(58,247)(663,249)
\Fill{14211B}(663,247)(1222,249)
\Fill{DFE3DF}(58,290)(663,291)
\Fill{DFE3DF}(663,290)(1222,291)
\Fill{DFE3DF}(58,333)(663,334)
\Fill{DFE3DF}(663,333)(1222,334)
\Fill{DFE3DF}(58,400)(663,401)
\Fill{DFE3DF}(663,400)(1222,401)
\Fill{DFE3DF}(58,443)(663,444)
\Fill{DFE3DF}(663,443)(1222,444)
\Fill{DFE3DF}(58,485)(663,486)
\Fill{DFE3DF}(663,485)(1222,486)
\Fill{14211B}(58,527)(663,529)
\Fill{14211B}(663,527)(1222,529)
% text
\T(58,77){\B{40}{0B5D4B}{Τι πήραμε από το ARGUS}}
\T(58,154){\R{20}{14211B}{Κάθε κίνδυνος, κάθε θέση αισθητήρα του ARGUS και κάθε συναγερμός στο παιχνίδι προέρχεται από τα ίδια τα δεδομένα του, και στα πέντε}}
\T(58,182){\R{20}{14211B}{μνημεία του.}}
\T(64,231){\B{19.5}{0B5D4B}{Από το ARGUS}}
\T(669.266,231){\B{19.5}{0B5D4B}{Στο AEGIS}}
\T(64,276){\R{19.5}{14211B}{Οι συντεταγμένες των εγκατεστημένων οργάνων}}
\T(669.266,276){\R{19.5}{14211B}{Η αποκάλυψη: τα όργανα του ARGUS στις θέσεις τους}}
\T(64,318){\R{19.5}{14211B}{Τα σενάρια κινδύνου και η καθημερινή αξιολόγησή τους}}
\T(669.266,318){\R{19.5}{14211B}{Οι τρεις κίνδυνοι κάθε αποστολής}}
\T(64,361){\R{19.5}{14211B}{Οι συναγερμοί Αυγούστου και Σεπτεμβρίου 2026, από όρια ειδικών και}}
\T(64,386){\R{19.5}{14211B}{από τεχνητή νοημοσύνη}}
\T(669.266,361){\R{19.5}{14211B}{Ο συναγερμός που ξαναζεί ο παίκτης}}
\T(64,428){\R{19.5}{14211B}{Τα φωτογραμμετρικά 3D μοντέλα και το έδαφος}}
\T(669.266,428){\R{19.5}{14211B}{Ο συναγερμός και η ζημιά σε 3D}}
\T(64,471){\R{19.5}{14211B}{Το API του ARGUS, μόνο για ανάγνωση, κάθε δύο ώρες}}
\T(669.266,471){\R{19.5}{14211B}{«Ζωντανά από το ARGUS» στην πρώτη οθόνη}}
\T(64,513){\R{19.5}{14211B}{Η μορφή εγγραφής πολιτιστικής τεκμηρίωσης (CIDOC CRM E14)}}
\T(669.266,513){\R{19.5}{14211B}{Η αναφορά παρατήρησης του επισκέπτη}}
\T(1157.297,701){\R{13}{88928C}{\LS{2}{AEGIS}}}
\T(1201.766,701){\R{13}{88928C}{\LS{2}{8 / 10}}}
\end{slide}

% ================================================= slide 9 — Τι κερδίζουν ο επισκέπτης και το ARGUS
\begin{slide}
\Fill{FBFCFA}(0,0)(1280,720)
\Fill{0B5D4B}(58,109)(1222,111)
\Shadow{images/slide09_shadow1.png}{0.1608}(791.04,130.24)(1232.96,388.8)[802,139,1222,376]
\Frame{D9DED9}{1}(802.5,139.5)(1221.5,375.5)
\Img{images/slide09_1.jpg}(803,140)(1221,375)
\Fill{A84300}(58,186)(69,197)
\Fill{A84300}(58,252)(69,263)
\Fill{A84300}(58,344)(69,355)
\Fill{A84300}(58,464)(69,475)
\Fill{A84300}(472,186)(483,197)
\Fill{A84300}(472,252)(483,263)
\Fill{A84300}(472,317)(483,328)
% text
\T(58,77){\B{40}{0B5D4B}{Τι κερδίζουν ο επισκέπτης και το ARGUS}}
\T(58,159){\B{21}{0B5D4B}{Ο επισκέπτης}}
\T(472.266,159){\B{21}{0B5D4B}{Το ARGUS}}
\T(802,399){\R{15}{5F6863}{Μια αναφορά παρατήρησης περιέχει πεδία όπως τι είδες, πού και}}
\T(802,421){\R{15}{5F6863}{πότε. Το παιχνίδι γίνεται δίαυλος επικοινωνίας, αλλά δεν στέλνει}}
\T(802,442){\R{15}{5F6863}{τίποτα μόνο του.}}
\T(84,197){\R{19.5}{14211B}{Κατανοεί τις ανάγκες και τους κινδύνους}}
\T(84,224){\R{19.5}{14211B}{κάθε μνημείου.}}
\T(84,263){\R{19.5}{14211B}{Φεύγει με έναν κώδικα επισκέπτη και μια}}
\T(84,290){\R{19.5}{14211B}{λίστα με όσα να προσέχει στην επόμενη}}
\T(84,316){\R{19.5}{14211B}{επίσκεψη.}}
\T(84,355){\R{19.5}{14211B}{Αν δει κάτι στο μνημείο που ίσως δεν έχει}}
\T(84,382){\R{19.5}{14211B}{προσέξει το ARGUS, το γράφει σε μια κάρτα}}
\T(84,409){\R{19.5}{14211B}{παρατήρησης και, αν θέλει, τη στέλνει στον}}
\T(84,436){\R{19.5}{14211B}{φορέα του μνημείου.}}
\T(84,475){\R{19.5}{14211B}{Ο εκπαιδευτικός βλέπει τις επιλογές της}}
\T(84,502){\R{19.5}{14211B}{τάξης δίπλα σε όσα εγκατέστησε το ARGUS.}}
\T(498.266,197){\R{19.5}{14211B}{Ένα δημόσιο πρόσωπο για τα}}
\T(498.266,224){\R{19.5}{14211B}{δεδομένα του.}}
\T(498.266,263){\R{19.5}{14211B}{Πολίτες εκπαιδευμένους πάνω σε}}
\T(498.266,290){\R{19.5}{14211B}{αυτά.}}
\T(498.266,328){\R{19.5}{14211B}{Παρατηρήσεις πολιτών ήδη στη}}
\T(498.266,355){\R{19.5}{14211B}{δική του μορφή εγγραφής,}}
\T(498.266,382){\R{19.5}{14211B}{έτοιμες αν αποφασίσει να τις}}
\T(498.266,409){\R{19.5}{14211B}{δεχτεί.}}
\T(1157.297,701){\R{13}{88928C}{\LS{2}{AEGIS}}}
\T(1201.766,701){\R{13}{88928C}{\LS{2}{9 / 10}}}
\end{slide}

% ================================================= slide 10 — Για όλους, και έτοιμο να το παίξετε
\begin{slide}
\Fill{FBFCFA}(0,0)(1280,720)
\Fill{0B5D4B}(58,109)(1222,111)
\Img{images/slide01_2.png}(58,371)(182,495)
\Shadow{images/slide10_shadow1.png}{0.1608}(951.04,132.16)(1218.88,683.84)[962,141,1208,671]
\Frame{D9DED9}{1}(962.5,141.5)(1207.5,670.5)
\Img{images/slide10_1.jpg}(963,142)(1207,670)
\Fill{A84300}(58,150)(69,161)
\Fill{A84300}(58,194)(69,205)
\Fill{A84300}(58,239)(69,250)
\Fill{A84300}(58,283)(69,294)
\Fill{A84300}(58,327)(69,338)
% text
\T(58,77){\B{40}{0B5D4B}{Για όλους, και έτοιμο να το παίξετε}}
\T(204,410){\B{21}{14211B}{Παίξτε το:}}
\T(204,439){\M{18}{14211B}{https://aloualou56.github.io/}}
\T(204,468){\M{18}{14211B}{argus\_trackA/}}
\T(84,163){\R{22}{14211B}{Πέντε γλώσσες: ελληνικά, αγγλικά, ιταλικά, ισπανικά, γερμανικά.}}
\T(84,207){\R{22}{14211B}{Ακουστική απόδοση κειμένων.}}
\T(84,252){\R{22}{14211B}{Σε υπολογιστή και κινητό, και σε σχολική τάξη.}}
\T(84,296){\R{22}{14211B}{Οδηγός στις πρώτες οθόνες, την πρώτη φορά.}}
\T(84,340){\R{22}{14211B}{Σχεδιασμένο με στόχο το επίπεδο AA του προτύπου προσβασιμότητας WCAG 2.2.}}
\T(1150.453,701){\R{13}{88928C}{\LS{2}{AEGIS}}}
\T(1194.922,701){\R{13}{88928C}{\LS{2}{10 / 10}}}
\end{slide}

\end{document}
